\chapter*{Declaration}
\addcontentsline{toc}{chapter}{Declaration}

I, \studentname{} (\rollnumber), student of \programme{} at \institution{}, declare that this report on the \emph{\projecttitle{}} is the result of my own work carried out under the guidance of \guidename{} during the academic year \academicyear{}. Sources and tools used are acknowledged in the report.

\vspace{2.5cm}
\hfill\begin{tabular}{c}
\rule{5cm}{0.4pt}\\
\studentname
\end{tabular}

\chapter*{Acknowledgement}
\addcontentsline{toc}{chapter}{Acknowledgement}

I thank \guidename{} for guidance and review throughout the project, the faculty and staff of \institution{} for sharing how student elections are run in practice, and my classmates who tested the system and gave feedback. I also thank the open-source communities whose tools made this work possible.

\chapter*{Abstract}
\addcontentsline{toc}{chapter}{Abstract}

Student elections in many colleges still rely on paper ballots, which are slow to count and hard to audit. This project delivers a self-hostable web platform that runs a campus election end to end. An installation wizard configures the institution; an Election Officer controls the election lifecycle; Desk Officers register students and issue Voter Slip IDs and single-use PINs; voters cast ballots on supervised booths that the officer has approved with a one-time token and an authenticator code. Ballots are stored apart from voter identity and duplicate votes are prevented inside a database transaction. Results stay hidden until the officer publishes them and can be revealed on a projector, one position at a time, or shared through a revocable public link. All important actions are recorded in an audit log. The system is packaged with Docker Compose and covered by 146 automated tests.

\textbf{Keywords:} electronic voting, campus election, secret ballot, audit log, Flask, React, PostgreSQL, Docker.

\tableofcontents
\listoffigures
\clearpage
